\section{Conclusion}

\begin{wrapfigure}{r}{0.53\linewidth}
    \centering
    \vspace{-0.5cm}

    % --- Left Image ---
    \begin{minipage}[t]{0.48\linewidth}
        \centering
        \begin{overpic}[width=\linewidth]{figures/limitation.pdf}
        \end{overpic}
        \caption{\textbf{Failure case 1.} Text prompt: \textit{``A swing hanging from a tree''}.
       }
        \label{fig:failure_case1}
    \end{minipage}
    \hfill
    % --- Right Image ---
    \begin{minipage}[t]{0.48\linewidth}
        \centering
        \begin{overpic}[width=\linewidth]{figures/fail2.png}
        \end{overpic}
        \caption{\textbf{Failure case 2.}
        Text prompt: \textit{``A brown leather sofa decorated with plush toys, \textbf{...}''}. See \autoref{supp:text_prompt} for the complete prompt.}
        \label{fig:failure_case2}
    \end{minipage}
    \vspace{-0.4cm}

\end{wrapfigure}
We presented \textbf{PAT3D}, a physics-augmented framework for text-to-3D scene generation that integrates vision-language reasoning with differentiable rigid body simulation. By decomposing the generation process into interpretable stages -- object and relation extraction, layout initialization, and physics-guided layout optimization -- our method produces 3D scenes that are not only semantically meaningful but also physically plausible and simulation-ready.
Through extensive experiments on diverse, contact-rich scenes, we demonstrated that PAT3D achieves superior physical realism compared to existing approaches. We believe PAT3D represents a step forward in bridging high-level scene understanding with low-level physical reasoning. We hope this work inspires further research in physically grounded, controllable, and editable 3D scene generation.


\paragraph{Limitations and Future Work}
We present two representative failure cases in \autoref{fig:failure_case1} and \autoref{fig:failure_case2}. While PAT3D handles most common physical dependencies, certain subtle relations remain challenging. In \autoref{fig:failure_case1}, the concept of “hanging” in the prompt “A swing hanging from a tree” is misinterpreted, whereas a physically correct configuration would require the swing to be suspended from two specific attachment points. Extending our system to a broader set of spatial relations and larger-scale scenes are both promising directions for future work. One potential avenue is to incorporate path planning techniques and adopt a hierarchical optimization strategy to manage spatial complexity more effectively.

Our simulation-in-the-loop optimization is solved using a local optimizer. Although the objective reliably decreases, it does not guarantee convergence to the global optimum corresponding to perfect semantic alignment. This limitation appears in the failure case of \autoref{fig:failure_case2}, where all stuffed toys should ideally rest on the sofa according to the prompt. However, due to the highly crowded configuration, the optimizer provides a suboptimal outcome: it reduces the number of fallen toys from three to one, yet one gray toy still remains on the floor. As the first work to incorporate differentiable simulation into 3D scene generation, we view the exploration of global optimization strategies as an exciting direction for future improvements, further expanding the possibilities opened by PAT3D.

